%&bericht

%%%%%%%%%%%%%%%%%%%%%%%%%%%%%%%%%%%%%%%%%%%%%%%%%%%%%%%%%%%%%%%%%%%%%%%%%%%%%%%
%% Descr:       Vorlage für Berichte der DHBW-Karlsruhe
%% Author:      Prof. Dr. Jürgen Vollmer, juergen.vollmer@dhbw-karlsruhe.de
%% $Id: bericht.tex,v 1.25 2020/03/13 15:07:45 vollmer Exp $
%%  -*- coding: utf-8 -*-
%%%%%%%%%%%%%%%%%%%%%%%%%%%%%%%%%%%%%%%%%%%%%%%%%%%%%%%%%%%%%%%%%%%%%%%%%%%%%%%

\documentclass[
   ngerman          % neue deutsche Rechtschreibung
  ,a4paper          % Papiergrösse
% ,twoside          % Zweiseitiger Druck (rechts/links)
% ,10pt             % Schriftgrösse
  ,11pt
% ,12pt
  ,pdftex
%  ,disable         % Todo-Markierungen auschalten
]{report}

% Bitte die Codierung Ihrer Dateien auswählen:
% \usepackage[latin1]{inputenc}    % Für UNIX mit ISO-LATIN-codierten Dateien
% \usepackage[applemac]{inputenc}  % Für Apple Mac
% \usepackage[ansinew]{inputenc}   % Für Microsoft Windows
\usepackage[utf8]{inputenc}        % UTF-8 codierte Dateien
                                   % Dieses Dokument ist unter Unix erstellt, daher
                                   % wird diese Input-Codierung benutzt.

\usepackage{bericht}
\usepackage{amsmath}
%% ACHTUNG, wenn man eine eigene Formatdatei (bericht.fmt) benutzt, werden Änderungen an bericht.sty
%% erst wirksam, wenn die Format-Datei neu erzeugt wurde!!!
%% Genauer alle Änderungen, die textuell vor der nächsten Zeile ".... endofdump...." stehen
%% werden erst wirksam, wenn die Formatdatei neu erzeugt wurde
\csname endofdump\endcsname

%%%%%%%%%%%%%%%%%%%%%%%%%%%%%%%%%%%%%%%%%%%%%%%%%%%%%%%%%%%%%%%%%%%%%%%%%%%%%%%
%% Angaben zur Arbeit
%%%%%%%%%%%%%%%%%%%%%%%%%%%%%%%%%%%%%%%%%%%%%%%%%%%%%%%%%%%%%%%%%%%%%%%%%%%%%%%

\newcommand{\Autor}{Yannic Hemmer}
\newcommand{\MatrikelNummer}{6853472}
\newcommand{\Kursbezeichnung}{TINF19B2}

\newcommand{\FirmenName}{SySS GmbH}
\newcommand{\FirmenStadt}{Tübingen}
%SELF \newcommand{\FirmenLogoDeckblatt}{\fbox{\includegraphics[width=3cm]{lion}}}

% Falls es kein Firmenlogo gibt:
%  \newcommand{\FirmenLogoDeckblatt}{}

\newcommand{\BetreuerFirma}{-}
\newcommand{\BetreuerDHBW}{Rolf Felder}

%%%%%%%%%%%%%%%%%%%%%%%%%%%%%%%%%%%%%%%%%%%%%%%%%%%%%%%%%%%%%%%%%%%%%%%%%%%%%%%%%%%%%

% Wird auf dem Deckblatt und in der Erklärung benutzt:
% \newcommand{\Was}{Projekt-/Studien-/Bachleorarbeit}
%\newcommand{\Was}{Projektrarbeit}
\newcommand{\Was}{Studienarbeit}
%\newcommand{\Was}{Bachleorarbeit}

%%%%%%%%%%%%%%%%%%%%%%%%%%%%%%%%%%%%%%%%%%%%%%%%%%%%%%%%%%%%%%%%%%%%%%%%%%%%%%%%%%%%%

\newcommand{\Titel}{Proof of Concept für den Angriff eines IT-Systems durch Implementierung kleptographischer Schwachstellen in kryptographischen Bibliotheken}
\newcommand{\AbgabeDatum}{16. Mai 2022}

\newcommand{\Dauer}{24 Wochen}

% \newcommand{\Abschluss}{Bachelor of Engineering}
\newcommand{\Abschluss}{Bachelor of Science}

\newcommand{\Studiengang}{Informatik / Informationstechnik}
% \newcommand{\Studiengang}{Informatik / Angewandte Informatik}

\hypersetup{%%
  pdfauthor={\Autor},
  pdftitle={\Titel},
  pdfsubject={\Was}
}

%%%%%%%%%%%%%%%%%%%%%%%%%%%%%%%%%%%%%%%%%%%%%%%%%%%%%%%%%%%%%%%%%%%%%%%%%%%%%%%

% Wenn \includeonly{..} benutzt wird, werden nur diese Kaptitel ausgegeben.
\includeonly{
  abk
 ,kapitel1
 ,kapitel2
 ,changelog
}

%%%%%%%%%%%%%%%%%%%%%%%%%%%%%%%%%%%%%%%%%%%%%%%%%%%%%%%%%%%%%%%%%%%%%%%%%%%%%%%

% Benutzt man das "biblatex"-Paket, dann muß das hier stehen:
% siehe auch die mit BIBLATEX markierten Zeilen in bericht.sty
\bibliography{bericht}

\begin{document}

%%%%%%%%%%%%%%%%%%%%%%%%%%%%%%%%%%%%%%%%%%%%%%%%%%%%%%%%%%%%%%%%%%%%%%%%%%%%%%%

\begin{titlepage}
\begin{center}
\vspace*{-2cm}
%SELF \FirmenLogoDeckblatt\hfill\includegraphics[width=4cm]{dhbw-logo}\\[2cm]
{\Huge \Titel}\\[1cm]
{\Huge\scshape \Was}\\[1cm]
{\large für die Prüfung zum}\\[0.5cm]
{\Large \Abschluss}\\[0.5cm]
{\large des Studienganges \Studiengang}\\[0.5cm]
{\large an der}\\[0.5cm]
{\large Dualen Hochschule Baden-Württemberg Karlsruhe}\\[0.5cm]
{\large von}\\[0.5cm]
{\large\bfseries \Autor}\\[1cm]
{\large Abgabedatum \AbgabeDatum}
\vfill
\end{center}
\begin{tabular}{l@{\hspace{2cm}}l}
Bearbeitungszeitraum	         & \Dauer 			\\
Matrikelnummer	                 & \MatrikelNummer		\\
Kurs			         & \Kursbezeichnung		\\
Ausbildungsfirma	         & \FirmenName			\\
			         & \FirmenStadt			\\
Betreuer der Ausbildungsfirma	 & \BetreuerFirma		\\
Gutachter der Studienakademie	 & \BetreuerDHBW		\\
\end{tabular}
\end{titlepage}

%%%%%%%%%%%%%%%%%%%%%%%%%%%%%%%%%%%%%%%%%%%%%%%%%%%%%%%%%%%%%%%%%%%%%%%%%%%%%%%

%%%%%%%%%%%%%%%%%%%%%%%%%%%%%%%%%%%%%%%%%%%%%%%%%%%%%%%%%%%%%%%%%%%%%%%%%%%%%%%
%% Descr:       Vorlage für Berichte der DHBW-Karlsruhe, Erklärung
%% Author:      Prof. Dr. Jürgen Vollmer, vollmer@dhbw-karlsruhe.de
%% $Id: erklaerung.tex,v 1.11 2020/03/13 14:24:42 vollmer Exp $
%% -*- coding: utf-8 -*-
%%%%%%%%%%%%%%%%%%%%%%%%%%%%%%%%%%%%%%%%%%%%%%%%%%%%%%%%%%%%%%%%%%%%%%%%%%%%%%%

% In Bachelorarbeiten muss eine schriftliche Erklärung abgegeben werden.
% Hierin bestätigen die Studierenden, dass die Bachelorarbeit, etc.
% selbständig verfasst und sämtliche Quellen und Hilfsmittel angegeben sind. Diese Erklärung
% bildet das zweite Blatt der Arbeit. Der Text dieser Erklärung muss auf einer separaten Seite
% wie unten angegeben lauten.

\newpage
\thispagestyle{empty}
\begin{framed}
\begin{center}
\Large\bfseries Erklärung
\end{center}
\medskip
\noindent
% siehe §5(3) der \enquote{Studien- und Prüfungsordnung DHBW Technik} vom 29.\,9.\,2017 und Anhang 1.1.13
Ich versichere hiermit, dass ich meine \Was mit dem Thema:
\enquote{\Titel}
selbstständig verfasst und keine anderen als die angegebenen Quellen und Hilfsmittel benutzt habe. Ich versichere zudem, dass die eingereichte elektronische Fassung mit der gedruckten Fassung übereinstimmt.
\vspace{3cm}
\noindent
\underline{\hspace{4cm}}\hfill\underline{\hspace{6cm}}\\
Ort~~~~~Datum\hfill Unterschrift\hspace{4cm}
\end{framed}

\vfill
\emph{Sofern  vom Dualen Partner ein Sperrvermerk gewünscht wird, ist folgende Formulierung
zu verwenden:}
\begin{framed}
\begin{center}
\Large\bfseries Sperrvermerk
\end{center}
\medskip
\noindent
Der Inhalt dieser Arbeit darf weder als Ganzes noch in Auszügen Personen
außerhalb des Prüfungsprozesses und des Evaluationsverfahrens zugänglich gemacht
werden, sofern keine anderslautende Genehmigung vom Dualen Partner vorliegt.
\end{framed}

%%%%%%%%%%%%%%%%%%%%%%%%%%%%%%%%%%%%%%%%%%%%%%%%%%%%%%%%%%%%%%%%%%%%%%%%%%%%%%%
\endinput
%%%%%%%%%%%%%%%%%%%%%%%%%%%%%%%%%%%%%%%%%%%%%%%%%%%%%%%%%%%%%%%%%%%%%%%%%%%%%%%

%%%%%%%%%%%%%%%%%%%%%%%%%%%%%%%%%%%%%%%%%%%%%%%%%%%%%%%%%%%%%%%%%%%%%%%%%%%%%%%

\begin{abstract}
Dieses \LaTeX-Dokument kann als Vorlage für einen Praxis- oder Projektbericht, eine Studien- oder
Bachelorarbeit dienen.

Zusammengestellt von Prof.\,Dr.\,Jürgen Vollmer \email{juergen.vollmer@dhbw-karlsruhe.de}\\
\url{https://www.karlsruhe.dhbw.de}. Die jeweils aktuellste Version dieses \LaTeX-Paketes ist immer
auf der \emph{FAQ-Seite} des Studiengangs Informatik zu finden:
\url{https://www.karlsruhe.dhbw.de/inf/studienverlauf-organisatorisches.html} $\to$ \emph{Formulare und Vorlagen}.

\centering Stand \verb+$Date: 2020/03/13 15:07:45 $+
\end{abstract}

\newpage
\tableofcontents           % Inhaltsverzeichnis hier ausgeben
\listoffigures             % Liste der Abbildungen
\listoftables              % Liste der Tabellen
\lstlistoflistings         % Liste der Listings
\listofequations           % Liste der Formeln

% Jetzt kommt der "eigentliche" Text
%%%%%%%%%%%%%%%%%%%%%%%%%%%%%%%%%%%%%%%%%%%%%%%%%%%%%%%%%%%%%%%%%%%%%%%%%%%%%%
%% Descr:       Vorlage für Berichte der DHBW-Karlsruhe, Datei mit Abkürzungen
%% Author:      Prof. Dr. Jürgen Vollmer, vollmer@dhbw-karlsruhe.de
%% $Id: abk.tex,v 1.4 2017/10/06 14:02:03 vollmer Exp $
%% -*- coding: utf-8 -*-
%%%%%%%%%%%%%%%%%%%%%%%%%%%%%%%%%%%%%%%%%%%%%%%%%%%%%%%%%%%%%%%%%%%%%%%%%%%%%%%

\chapter*{Abkürzungsverzeichnis}                   % chapter*{..} -->   keine Nummer, kein "Kapitel"
						         % Nicht ins Inhaltsverzeichnis
% \addcontentsline{toc}{chapter}{Akürzungsverzeichnis}   % Damit das doch ins Inhaltsverzeichnis kommt

% Hier werden die Abkürzungen definiert
\begin{acronym}[DHBW]
  % \acro{Name}{Darstellung der Abkürzung}{Langform der Abkürzung}
 \acro{Abk}[Abk.]{Abkürzung}

 % Folgendes benutzen, wenn der Plural einer Abk. benöigt wird
 % \newacroplural{Name}{Darstellung der Abkürzung}{Langform der Abkürzung}
 \newacroplural{Abk}[Abk-en]{Abkürzungen}


 \acro{RSA}[RSA]{Rivest-Shamir-Adleman}
 \acro{AES}[AES]{Advanced Encryption Standard}

\end{acronym}
              % Abkürzungsverzeichnis

\chapter{Grundlagen}

In diesem Kapitel werden die theoretischen Grundlagen der Kryptologie, Mathematik, Zahlentheorie und der IT-Sicherheit erläutert, die in dieser Arbeit eine Rolle spielen. Aus diesen sollen die grundlegenden Funktionen eines kleptographischen Angriffes und dessen Folgen abgeleitet werden.

\section{Kryptologie}
    Die Kryptologie ist die wissenschaftliche Disziplin für den Schutz von Daten. Unter ihr stehen die zwei Felder der Kryptographie und der Kryptoanalyse.

    \subsection{Kryptographie}
        Die Kryptologie befasst sich mit der Entwicklung von Verfahren und Techniken für den sicheren Austausch von Daten. Dabei stehen zwei Eigenschaften \cite{BSW.2015} im Fokus:

        \subsubsection{Eigenschaften}
            \paragraph{Geheimhaltung}
                Durch Geheimhaltung sollen, bei der Übertragung von Daten zwischen Teilnehmern, Unbeteiligte keine Erkenntnisse über den Inhalt erlangen. Dies kann durch physikalische oder organisatorische Maßnahmen erreicht werden, wobei Unbeteiligten der Zugang zu den übertragenen Daten verwehrt wird. Diese Maßnahmen sind sinnvoll bei der Übergabe der Daten in einer nicht digitalen Welt. Bei der Kommunikation in digitalen Netzen, wie u.a. dem Internet, sind diese Maßnahmen nur schwer zu implementieren. Dies gilt nicht für kryptographische Maßnahmen. Dabei ist es nicht mehr das Ziel, Unbeteiligten den Zugang zu den übertragenen Daten zu erschweren, sondern den Inhalt der Daten während der Übertragung zu verschlüsseln. Dadurch soll es Unbeteiligten nahezu unmöglich sein, aus den mitgehörten oder abgefangenen Daten, Rückschlüsse auf deren Inhalt zu erlangen.
                
            \paragraph{Authentifikation}
                Durch Authentifikation soll es den Teilnehmern einer Kommunikation möglich sein, die anderen Teilnehmer und empfangene Nachrichten zweifelsfrei identifizieren und zuweisen zu können. Hierbei spielen Signaturverfahren eine wichtige Rolle, da kein Geheimnis benötigt wird um einen Teilnehmer zu authentifizieren. Dabei können sich Teilnehmer durch das Wissen oder den Besitz eines Geheimnisses (Passwort, Zertifikat, Schlüssel) authentifizieren. \cite{BSW.2015}

        Nur wenn beide Eigenschaften gegeben sind ist eine Übertragung von Daten als sicher anzusehen. Falls die Geheimhaltung fehlt, kann der Inhalt durch Sniffing mitgelesen werden. Falls die Authentifikation der Teilnehmer fehlt, können sich Unbeteiligte als "echte" Teilnehmer ausgeben und somit die Daten an ihrem Endpunkt entschlüsseln.

        \subsubsection{Verschlüsselungsverfahren}
        % Was sind Verschlüsselungsverfahren, Ziel, Historie, momentaner Stand, State of the Art (AES, RSA )
            \paragraph{Asymmetrische Verschlüsselung}
            % Hier angesprochene Verfahren, unterschiede, Definition / Abgrenzung zu symmetrischen Verfahren

            \paragraph{Hybride Verfahren}
                % Meiste Verfahren heute hybride Verfahren. Asymmetrisch zum Schlüsselaustausch. Symmetrisch zum Verschlüsseln. Laufzeit von Asymmetrischen Verfahren. Wenn Schlüsselaustausch sicher Geheimhaltung verletzt.

        \subsubsection{Angriffe}
            % Bedeutung von Brute-Force bei modernen kryptographischen Verfahren

            % Ziel der Angriffe / Bedeutung von Schlüsseln
            
            \paragraph{Known Cipher Attack}

            \paragraph{Known Plaintext Attack}

            \paragraph{Chosen Plaintext Attack}

            \paragraph{Chosen Cipher Attack}


\section{Mathematik}
    Mathematische Probleme stellen die Grundlage für moderne Kryptographie. 

    \subsection{Diskreter Logarithmus}
        Bei der Bestimmung des Logarithmus wird der Exponent (hier: $x$) gesucht, welcher mit einer bekannten Zahl als Basis $z$, eine weitere bekannte Zahl $y$ ergibt.
        \begin{equation}
            z^{x} = y
            \eqlabel{eq-log}{Normaler Logarithmus}
        \end{equation}

        Der diskrete Logarithmus bezieht hier auf die Berechnung des Logarithmus in ein Gruppe. Diese Gruppe bildet sich aus der Rechnung mit Restklassen (modulo). Dadurch entsteht folgendes Problem, bei der die Variable $x$ gesucht ist und alle anderen Variablen bekannt sind.
        \begin{equation}
            z^{x} \pmod n \equiv y
            \eqlabel{eq-dis-log}{Diskreter Logarithmus}
        \end{equation}
        
        Hierbei ist in der Notation zu beachten, dass sich durch das Rechnen auf mit einer Gruppe, Äquivalenzklassen ($\equiv$) bilden. Diese entsprechen den Restklassen des Rechnen mit Modulo. $n$ ist die Mächtigkeit der Aquivalenzklassen.

        Die Bestimmung von $x$ in \ref{eq-dis-log} wird als Problem des diskreten Logarithmus bezeichnet. Mit der Komplexität wird sich in den Grundlagen der Komplexitätstheorie beschäftigt.

        Dabei ist die Umkehrfunktion, des diskreten Logarithmus $f(x)$ \ref{eq-dis-log}, mathematisch einfach zu berechnen. Diese Umkehrfunktion entspricht der diskreten Exponentialfunktion:
        \begin{equation}
            f^{-1}(x) = z^{x} \pmod n \equiv y
            \eqlabel{eq-dis-exp}{Diskretere Exponentialfunktion}
        \end{equation}
        Hierbei sind $z,x,n$ gegeben und $y$ gesucht.

    \subsection{Faktorisierung}


    \subsection{Effiziente Berechnung der diskreten Exponentialfunktion}
        In der Kryptographie werden große Zahlen genutzt, um die Sicherheit der verwendeten Algorithmen zu gewährleisten. Hierfür wird als Beispiel angenommen, dass als Basis $z$ eine 256-bit Lange Zahl hoch einem 300-bit langem Exponenten $x$ genommen werden soll. Hierbei ist $n$ 1024-bit lang. 

        Wenn mann nun $z$ in Byte berechnet wäre dies eine 32 Byte lange Zahl.

        $x$ entspricht einer ungefähr 90. stelligen Zahl. 

        \begin{equation}
            z^{10^{90}} \pmod n \equiv y
            \eqlabel{eq-dis-exp-lN}{Diskretere Exponentialfunktion mit großen Zahlen}
        \end{equation}

        Eine numerische Berechnung von $ z^{10^{90}} $ ist aufgrund von begrenzten Ressourcen nicht möglich. 

        Jedoch kann man sich die diskrete Eigenschaft dieser Problems sich zu nutzte machen. Hierfür können Verfahren, wie Square-and-Multiply zusammen mit der Restklassenberechnung genutzt werden. Dadurch lassen sich auch großzahlige Exponenten berechnen. Hierfür soll ein einfaches Beispiel gegeben werden:
        \begin{equation}
            37^{52} \pmod {128} \equiv y
            \eqlabel{eq-dis-exp-bspOne}{Diskretere Exponentialfunktion mit großen Zahlen Beispiel-Eins}
        \end{equation}
        Bei Betrachtung der Äquivalenzgleichung fällt auf, dass $37^{52}$ eine große Zahl ergibt. Jedoch wird diese Zahl noch $ x \pmod 128$ gerechnet. Dadurch liegt das Ergebnis in einem Zahlenraum von:
        \begin{equation}
            y \in \mathbb{N} \mid 0 \le y < 128 
            \eqlabel{eq-dis-exp-numSpace}{Diskretere Exponentialfunktion in Zahlenraum}
        \end{equation}

        Auf Grundlage der Potenzgesetze wird $37^{52}$ nun zerlegt. 
        \begin{equation}
        \begin{aligned}
            52 = 32 + 16 + 4 = 2^{5} + 2^{4} + 2^{2} \\
            37^{52} \pmod {128} \equiv 37^{2^{5}} * 37^{2^{4}} * 37^{2^{2}} \\
            \equiv 37^{2^{5}} \pmod {128} * 37^{2^{4}} \pmod {128} * 37^{2^{2}} \pmod {128}
        \end{aligned}
        \eqlabel{eq-dis-exp-bspTwo}{Diskretere Exponentialfunktion mit großen Zahlen Beispiel-Zwei}
        \end{equation}

        Die einzelnen Bestandteile werden dann iterativ berechnet und durch Multiplikation zusammengefasst (siehe \ref{eq-dis-exp-bspTwo}). Dies wird als Square-and-Multiply-Verfahren bezeichnet.
        \begin{equation}
        \begin{aligned}
            37^{2^{2}} \pmod {128} \equiv (37^{2^{1}} \pmod {128})^{2} \\
            37^{2^{3}} \pmod {128} \equiv (37^{2^{2}} \pmod {128})^{2} \\
            37^{2^{4}} \pmod {128} \equiv (37^{2^{3}} \pmod {128})^{2} \\
            37^{2^{5}} \pmod {128} \equiv (37^{2^{4}} \pmod {128})^{2}
        \end{aligned}
        \eqlabel{eq-dis-exp-bspThree}{Diskretere Exponentialfunktion mit großen Zahlen Beispiel-Drei}
        \end{equation}
        
        \subsubsection{Allgemein}
        Gegeben mit gesucht $y$:
        \begin{equation}
            z^{x} \pmod n \equiv y
        \end{equation}

        Zerlegung von $x$ eine Summe von Zweierpotenzen:
        \begin{equation}
            x = 2^{0} + 2^{1} + 2^{2} + ...
        \end{equation}

        Dabei bilden die binären Logarithmen der einzelnen Zweierpotenzen die Menge $\mathbb{K}$.

        Berechnung der einzelnen Faktoren durch iteratives Square-and-Multiply-Verfahren. Dies wird bis $f(max(\mathbb{K}))$ berechnet. $max(\mathbb{K})$ steht hier für das Element von $\mathbb{K}$, mit dem größten Wert.
        \begin{equation}
            f(i+1) = f(i)^{2} \pmod n \mid f(1) = z^{1} \pmod n 
        \end{equation}

        Zuletzt wird das Produkt, aller Ergebnisse von $f(x)$ für die Elemente der Menge $\mathbb{K}$, gebildet. Dabei gilt:

        \begin{equation}
            \prod_{k \in \mathbb{K}} f(k) \equiv z^{x} \pmod n \equiv y
        \end{equation}





\chapter{RSA}
\ac{RSA} ist ein kryptographisches Verfahren, welches zu den Public-Key-Verfahren gehört. Der Verfahren wurde von R. Rivest, A. Shamir und L. Adleman entwickelt und trägt deshalb ein Anagram der Erfinder als Namen.

\section{Ablauf} \label{rsa-ablauf}
    \paragraph{Ausgangszenario} Teilnehmer A will über ein öffentliches Netz sicher mit anderen Teilnehmer kommunizieren. 

    \paragraph{Schlüsselgeneration} Damit andere Teilnehmer geheime Nachrichten schicken können muss A sich ein Schlüsselpaar generieren. Dafür wählt er zwei zufällige und große Primzahlen: $P$ und $Q$. 

    Das Produkt von $P$ und $Q$ bildet $N$, welche den Modulo / den Zahlenraum für weitere mathematische Operationen festlegt: 
        \begin{equation}
            N = P \cdot Q
            \eqlabel{eq-rsa-npq}{RSA - Primfaktoren}
        \end{equation}

    Daraufhin berechnet der Teilnehmer die Eulersche $\phi$-Funktion von $P$ und $Q$:
        \begin{equation}
            \phi(N) = (P-1) \cdot (Q-1)
            \eqlabel{eq-rsa-phi_n}{RSA - Eulersche $\phi$-Funktion}
        \end{equation})
    
    Der öffentliche Schlüssel $E$ ist dann eine zu $\phi(N)$ teilerfremde Zahl. Man kann dies auch vereinfachen und eine Fermat'sche Primzahl für $E$ verwenden:
        \begin{equation}
            2^{2^{N}}+1 \mid N \in \{0,1,2,3,4\}
            \eqlabel{eq-rsa-ferma}{Fermat'sche Primzahl}
        \end{equation}
    
    Der größte gemeinsame Teiler von $E$ und $\phi(n)$ ergibt $1$, wobei sich der private Schlüssel $D$ aus der Vielfachsummendarstellung ergibt. Die Vielfachsummendarstellung wird dem euklidischen Algorithmus entnommen.  Dabei gilt:
        \begin{equation}
            c \cdot \phi(N) + E \cdot D \equiv 1
            \eqlabel{eq-rsa-calc_d}{RSA - Berechnung des privaten Schlüssels}
        \end{equation}
    
    Danach besitzt der Teilnehmer A ein öffentlichen Schlüssel $E$ und einen privaten Schlüssel $D$. Er kann nun den öffentlichen Schlüssel $E$ zusammen mit $N$ veröffentlichen. $E$ zusammen mit $N$ und $D$ sind ein Schlüsselpaar welches in einen öffentlichen Teil und ein privaten Teil geteilt wird. Dabei herrscht eine eindeutige Zuordnung zwischen den Teilen.
    
    \subsection{Verschlüsselung} \label{sec-RSA-crypt}
        Wenn Teilnehmer B eine verschlüsselte Nachricht $M$ an Teilnehmer A senden will, braucht er hierfür den öffentlichen Schlüssel von A. $M$ muss aber dabei kleiner als $N$ sein. Dabei B  verschlüsselt wie folgt, wobei $C$ der resultierende Geheimtext ist.

        \begin{equation}
            M^{E_{A}} \pmod N = C
            \eqlabel{eq-rsa-encrypt}{RSA - Verschlüsselung einer Nachricht}
        \end{equation}

        Wenn Teilnehmer A den Geheimtext $C$ von B erhält, entschlüsselt er diesen mit seinem privaten Schlüssel. Dadurch berechnet er die von Teilnehmer B verschlüsselte Nachricht $M$.

        \begin{equation}
            C^{D_{A}} \pmod N = M
            \eqlabel{eq-rsa-decrypt}{RSA - Entschlüsselung eines Geheimtextes}
        \end{equation}

        Da $D$ privat und eindeutig für den öffentlichen Teil des Schlüsselpaars ist, können nur Teilnehmer, die über $D$ verfügen, Nachrichten entschlüsseln, die mit dem zugehörigen öffentlichen Teil verschlüsselt wurden.
    
    \subsection{Signatur} \label{sec-RSA-sign}
        Das RSA-Verfahren ermöglicht auch das Signieren von Nachrichten. Dies wird z.B. im Internet genutzt um Teilnehmer zu authentifizieren und Nachrichtenintegrität zu sichern. RSA wird jedoch meist nur auf kleinere Nachrichten wie Fingerabdrücke angewandt.

        Signatur $S$ einer Nachricht $M$ von Teilnehmer A:
        \begin{equation}
            M^{D_{A}} \pmod N = S
            \eqlabel{eq-rsa-sign}{RSA - Signieren einer Nachricht}
        \end{equation}

        Zum Überprüfen der Echtheit braucht der Teilnehmer B den öffentlichen Schlüssel von A:
        \begin{equation}
            S^{E_{A}} \pmod N = M
            \eqlabel{eq-rsa-check_sign}{RSA - Prüfen einer Signatur}
        \end{equation}

        Eine Signatur kann eindeutig Teilnehmer A zugeordnet werden, da nur ein Teilnehmer im Besitz von $d_{A}$ eine Signatur einer Nachricht erstellen kann, die mit dem öffentlichen Teil des Schlüsselpaars von A überprüft werden kann.
    
\section{Annahmen} \label{sec-RSA-assertion}
        Das Ergebnis einer \ac{RSA}-Verschlüsselung wird als Zufallszahl betrachtet. Grundlage hierfür ist, dass RSA mit sicheren Schlüsseln und Implementation, einer Einweg-Funktion gleicht, wenn der private Schlüssel nicht bekannt ist. Das Ergebnis einer \ac{RSA}-Verschlüsselung ist folgendermaßen die Generierung einer Zufallszahl $C$:
        \begin{equation}
            C \in [1,N)
            \eqlabel{eq-rsa-random}{RSA-Verschlüsselung EInwegfunktion}
        \end{equation}
        % Kann C 0 oder 1 sein? 
        Diese Annahme ist für die mathematische Betrachtung von möglichen Grenzfälle später relevant. Im kryptographischen Kontext, kann natürlich die Einwegfunktion mit Kenntnis des privaten Schlüssels effizient umgekehrt werden.

\section{Sicherheit}
    Die Sicherheit des RSA-Verfahrens basiert auf zwei mathematischen Problemen, welche unter Aufwand endlicher Ressourcen, nicht gelöst werden können. Hierbei wird sich sowohl auf RSA-gestützte Verschlüsselungs- und Signaturverfahren bezogen.
    Diese zwei Probleme sind:
    \begin{itemize}
        \item Faktorisierung einer bekannten Zahl, welche das Produkt zweier großer Primzahlen ist. Im Kontext von RSA ist diese Zahl mit $N$ repräsentiert.
        \item Bestimmung des diskreten Logarithmus. Bei RSA wäre dies die Bestimmung von 
        \begin{equation}
            D \mid M^{D} \equiv C \pmod N .
        \end{equation}
    \end{itemize}

    Für die Sicherheit der Public-Key-Verschlüsselung von RSA, spielt nicht effiziente Berechenbarkeit der Faktorisierung die Hauptrolle. Falls mit RSA signiert werden soll, ist zusätzlich die Unberechenbarkeit des diskreten Logarithmus wichtig. Ansonsten könnte der private Schlüssel bestimmt werden.

    \paragraph{Angriffe auf RSA}
    % Momentane Angriffe auf RSA
    % RSA ist momentan als sicher bekannt

    % Fermatsche Faktorisierung
    % Gemeinsame Primzahlen

\chapter{Kleptographie}
    \section{Definition}
    % Was ist Kleptographie


    \section{Geschichte}
    % Ursprünglicher Angriff / Konzept

    \section{Angriffskategorie}
        Bisher wurden Angriffe auf kryptographisches Systeme in eine der vier Kategorien \ref{sec-Kryptoanalyse} unterteilt (Known Cipher, Known Plaintext, Chosen Cipher, Chosen Plaintext). Ein kleptographischer Angriff fällt jedoch in keiner dieser Kategorien. Für kleptographische Angriffe müsste eine weitere, fünfte Kategorie geschaffen werden: Known Key Attacks. 
        % Angriff auf die Implementation

    \section{Aufbau kleptographischer Angriffe}
    

        \subsection{Voraussetzungen} 
        Die Voraussetzungen für den hier beschriebenen kleptographischen Angriff ist die einmalige Manipulation eines kryptographischen Systems zur Erzeugung von \ac{RSA}-Schlüsselpaaren. Dies kann während der Herstellung oder des Betriebs des Systems gelingen. Dabei muss nur die Schlüsselerstellung manipuliert werden. Die Funktionen zum Signieren, Verifizieren, Ver- und Entschlüsseln des Systems sind davon sich betroffen.


        \subsection{SETUP}
        % Beschreibung des SETUPS

    \section{SETUP für RSA}
        \subsection{Voraussetzungen}
            Für ein \ac{SETUP}-Angriff auf eine Implementation von \ac{RSA} hat der Angreifer ein eigenes Schlüsselpaar: $N_{A}$ Modulus des Angreifers, $E_{A}$ Öffentlicher Schlüssel des Angreifers, $D_{A}$ Privater Schlüssel des Angreifers. Das Schlüsselpaar wird wie für \ac{RSA} üblich generiert.
        
        \subsection{Generierung und Verschlüsselung}
            \paragraph{Schritt 1} \label{sec-Schritt-Gen 1} Es wird eine zufällige Primzahl $P$ generiert. $P$ wird dann mit dem öffentlichen Schlüssel des Angreifers verschlüsselt:
            \begin{equation}
                vP = P^{E_{A}} \mod N_{A}
                \eqlabel{eq-SETUP-vP}{Verschlüsselung von P mit dem öffentlichen Schlüssel des Angreifers}
            \end{equation}

            \paragraph{Schritt 2} \label{sec-Schritt-Gen 2} $N'$ wird gebildet indem $vP$ und eine Zufallszahl gleicher Länge $t$ in binärer Form konkateniert werden:
            \begin{equation}
                N' = vP || t
                \eqlabel{eq-SETUP-N'}{Berechnung des temporären Modulus}
            \end{equation}
            $N'$ ist dabei nicht der Modulus des generierten \ac{RSA}-Schlüsselpaars sondern nur eine temporäre Form.

            \paragraph{Schritt 3} \label{sec-Schritt-Gen 3} Berechnung der zweiten Primzahl $Q$: 
            \begin{equation}
                P \cdot Q + R = N'
                \eqlabel{eq-SETUP-Q}{Berechnung der zweiten Primzahl P}
            \end{equation}
            
            \paragraph{Schritt 4} \label{sec-Schritt-Gen 4} Bestimmung des Modulus $N$, wie für \ac{RSA} üblich durch:
            \begin{equation}
                N = P \cdot Q
                \eqlabel{eq-SETUP-N}{Berechnung von N}
            \end{equation}

            \paragraph{Schritt 5} \label{sec-Schritt-Gen 5} Wählen des öffentlichen Schlüssels $E$ und Berechnen des privaten Schlüssels $D$ mittels der modularen multiplikativen Inversen bzgl. $\phi(N)$:
            \begin{equation}
                D = modular\_multiplicative\_inverse(E, \phi(N))
                \eqlabel{eq-SETUP-D}{Berechnung von D}
            \end{equation}

            \paragraph{Schritt 6} \label{sec-Schritt-Gen 6} Mit Schritt 5 wurde ein vollkommen funktionales \ac{RSA}-Schlüsselpaar erstellt. Mittels diesem können nun Informationen verschlüsselt/signiert, Chiffren entschlüsselt und Signaturen verifiziert werden, wie in \ref{sec-RSA-crypt} und \ref{sec-RSA-sign} gezeigt wurde.
        
        \subsection{Angriff}
            \paragraph{Schritt 1} \label{sec-Schritt-Ang 1} Der Angreifer erlangt den öffentlichen Schlüssel des Ziels und besitzt somit $N$ und $E$. Dies ist möglich, da diese Informationen öffentlich sind.

            \paragraph{Schritt 2} \label{sec-Schritt-Ang 2} Der Angreifer teilt $N$ in binärer Form in der Hälfte womit er $vP$ erhält. Die mathematische Begründung hierfür in \ref{sec-SETUP-vP_from_N}.

            \paragraph{Schritt 3} \label{sec-Schritt-Ang 3} $P$ wird durch die Entschlüsslung von $vP$ mittels des privaten Schlüssels des Angreifers berechnet: 
            \begin{equation}
                P = (vP)^{D_{A}} \mod N_{A}
                \eqlabel{eq-SETUP-P}{Berechnung des ersten Primfaktors bei kleinem R}
            \end{equation}
            Damit ist dieser Schritt die inverse Operation zu \ref{sec-Schritt-Gen 1}.
            Zusätzlich muss auch $vP + 1$ entschlüsselt werden.
            Die mathematische Begründung hierfür in \ref{sec-SETUP-vP_from_N}.
            % Bei der Überprüfung, Prüfen ob P und Q prim
            \begin{equation}
                P' = (vP+1)^{D_{A}} \mod N_{A}
                \eqlabel{eq-SETUP-Palt}{Berechnung des ersten Primfaktors bei großem R}
            \end{equation}
            
            \paragraph{Schritt 4} \label{sec-Schritt-Ang 4} Hiermit ist der Angreifer im Besitz des ersten Primfaktors $P$ oder $P'$. Somit ist die Primfaktorzerlegung von $N$ trivial:
            \begin{equation}
                Q = N / P
                \eqlabel{eq-SETUP-Q}{Primfaktorzerlegung für P}
            \end{equation}
            Die Primfaktorzerlegung muss, gleich wie bei \ref{sec-Schritt-Ang 3}, für den alternativen Primfaktor $P'$ berechnet werden:
            \begin{equation}
                Q' = N / P'
                \eqlabel{eq-SETUP-Q}{Primfaktorzerlegung für P'}
            \end{equation}

            \paragraph{Schritt 5} \label{sec-Schritt-Ang 5} Der Angreifer ist hiermit im Besitz der Primfaktoren $P$, $Q$ und kann den privaten Schlüssel $D$ bestimmen \eqref{eq-SETUP-D}. Gleiches muss für die alternativen Primfaktoren berechnet werden.

            \paragraph{Schritt 6} \label{sec-Schritt-Ang 6} Der Angreifer besitzt den privaten und öffentlichen Schlüssel. Somit können Chiffren entschlüsselt und Signaturen gefälscht werden. Dabei muss der Angreifer, wenn noch nicht geschehen, den privaten Schlüssel $D$ und den alternativen privaten Schlüssel $D'$ einmalig testen, um den richtigen zu bestimmen.

        \subsection{Hintergründe zum RSA-SETUP}
            \subsubsection{Informationsgewinnung von vP aus N} \label{sec-SETUP-vP_from_N}
            Längendefinitionen:
            \begin{itemize}
                \item $N, N'$ hat $n$ Bit
                \item $vp, R$ hat $n/2$ Bit
            \end{itemize}
            $vP$ kann aus $N$ bestimmt werden, durch die Beziehung von $N$ und $N'$. Dabei gilt $N + R = N'$, wobei $R$ dem Angreifer nicht bekannt ist. Jedoch ist definiert, dass $R$ nur halb so viele Bits wie $N$ hat. Dadurch findet die Addition von $N + R$ nur auf den niederwertigeren Bits von $N$ statt. Somit sind die höherwertigen Bits von $N$ (Bits > $(n/2)$) nicht von der Addition mit $R$ beeinflusst und somit gleich zu den höherwertigen Bits von $N'$. Dies ist nur dann nicht der Fall, wenn es bei der Addition zu einem Überlauf für. In diesem Fall sind alle Bit > $(n/2)$ von $N$, die höherwertigen Bits von $N' + 1$.
            Zudem ist definiert, dass die höherwertigen von $N'$ gleich $vP$ sind (siehe \ref{eq-SETUP-N'}).
            Da der Angreifer nur $N$ aber nicht auch $R$ kennt, muss die Möglichkeit eines Überlauf durch $R$ berücksichtigt werden. 
            Aufgrund dessen kann der Angreifer $vP$ aus den höherwertigen Bits von $N$ lesen. Muss aber auch $vP + 1$ als Option berücksichtigen, da er nicht weiß, ob ein Überlauf stattgefunden hat.               

            \subsection{Verfahren zur Bestimmung der korrekten Primfaktoren}
                Die Schritte 3 bis 6 des Angriffs befassen sich mit dem Finden des korrekten Primfaktor aus den zwei resultierenden Möglichkeiten von vP $vP$ und $vP+1$. Daraus werden die Werte und ihre Alternativen für $P$, $N$, $Q$ und $D$ berechnet. 
                Um schlussendlich zu entscheiden, ob die Werte die aus $vP$ oder $vP+1$, kann eine Signatur mit $D$ und $D'$ mit einer Signatur des Angriffsziels verglichen werden. Dadurch kann eine eindeutige Entscheidung getroffen werden. 

                Diese Entscheidung kann jedoch unter Umständen früher berechnet werden. Dieser Fall kann bei folgenden Berechnungsschritten  auftreten:

                \subsubsection{Berechnung von P}
                    $P$ wird berechnet indem $vP$ mittels dem privaten Schlüssel $D_{A}$ des Angreifers entschlüsselt wird. Dabei sollte, wie für eine \ac{RSA}-Ver-/Entschlüsselung üblich, eine vollkommen zufällige Zahl resultieren. Jedoch ist es eine Bedingung, dass $P$ prim ist. 
                    Die Wahrscheinlichkeit, dass eine Zufallszahl, mit steigender Anzahl an Stellen, prim ist sehr gering. 
                    Mathematische Erläuterung %ref

                    Falls das entschlüsselte $P$ nicht prim ist, muss es die Alternative $P'$ sein. Gleiches gilt wiederum auch für $P'$.

                \subsubsection{Berechnung von Q} 
                    Bei der Berechnung von $Q$ gilt die gleiche Eigenschaft, wie bei $P$, dass $Q$ prim seien muss.

                    Jedoch kann bei der Berechnung von $Q$ eine eindeutige Entscheidung getroffen werden. Der Grund dafür ist, dass $Q$ eine ganze Zahl seien muss, also $Q \in \mathbb(Z)$. 
                    Ohne weitere geltende Bedingungen wäre die Überprüfung von $Q = N/P$ und $Q' = N/P'$ auf 
                    \begin{equation}
                        \text{Correct prime factors} =
                        \begin{cases}
                            (P, Q) & Q \in Z \wedge Q' \notin \mathbb(Z) \\
                            (P', Q') & Q' \in Z \wedge Q \notin \mathbb(Z) \\
                            (P, Q) & sonst
                        \end{cases}
                        \eqlabel{eq-OPT-Q}{Optimierung für die Bestimmung von Q}
                    \end{equation}
                    Ohne geltende Bedingungen von \ac{RSA} könnte $Q, Q' \in \mathbb(Z)$ gelten. Da $N$ ein vielfaches beider Zahlen $P$ und $P'$ seien kann. Jedoch ist $N$ in \ac{RSA} das Produkt von zwei Primzahlen. Dadurch gibt es zwei Zahlen $x$ für die gilt $ (N/x) \in \mathbb(Z) $. Dabei handelt es sich um die korrekten Primfaktoren $P$ und $Q$. Dadurch gilt immer nur einer der beiden Fälle in \ref{eq-OPT-Q}. 
                    Im Fall 1 ist $ (N/P) \in \mathbb{Z}$, während $N$ kein vielfaches von $P'$ ist.
                    Umgekehrtes gilt für Fall 2.
                    Der dritte Fall wird dennoch benötigt. Er tritt ein, wenn $Q' = P$ oder $Q = P'$ gilt. 
                    Dieser Fall tritt dann ein, wenn $(P^{E} + 1)^{D} \mod N = Q$ gilt. Unter der Annahme, dass die Entschlüsselung einer Hash-Funktion gleich, tritt dieser Fall mit einer Wahrscheinlichkeit von $p(1/N)$ auf. 
                    Dieser sehr unwahrscheinliche Fall wird durch Fall 3 von \ref{eq-OPT-Q} abgedeckt. In diesem Fall wird eine der beiden Möglichkeiten ausgewählt, da durch die Kommutativität der Multiplikation es egal ist, ob $[P, Q] \vee [Q, P]$ die richtigen Primfaktoren von $N$ sind.

                    Durch \ref{eq-OPT-Q} kann sehr einfach und effizient die richtigen Primfaktoren ausgewählt werden. Diese Überprüfung hat keinen Einfluss auf die Laufzeit des Algorithmus.

                    Somit ist diese Überprüfung nicht nur eindeutiger als eine Überprüfung von $P$ auf prim, sondern auch effizienter.

                \subsubsection{Fehler bei der modularen multiplikativen Inverse}
                    Bei der Berechnung von $D$ wird die modulare multiplikative Inverse von $E$ und $\phi(N)$ bestimmt. Dabei kann es zu einem Fehler kommen, da für den Tupel von $E$ und $\phi(N)$ möglicherweise keine modulare multiplikative Inverse existiert.

                    % Mathematische Wahrscheinlichkeit
                
                




        



      



\chapter{Angriffskonzept}

    \section{Ziel}
        Folgender Abschnitt der Arbeit, befasst sich mit der Auswahl des Ziels für einen kleptographischen Angriff. Hierbei soll es sich um eine Softwarebibliothek handeln. Zudem soll diese Open-Source und hinreichend verbreitet sein. Die zugrundeliegende Programmiersprache soll abstrakt genug sein, dass die kryptographischen Operationen in Software abgebildet werden. Programmiersprachen, welche Hardware, wie \ac{TPM} nutzten, sind nicht für einen Angriff auf Softwarebibliotheken weniger geeignet.

        Als Programmiersprache wurde für die Implementation Python gewählt, weil Python weit verbreitet ist, eine Vielzahl an Open Source Libraries und die Syntax nah an Pseudocode liegt. Somit kann die Implementation leichter verstanden und schneller in andere Sprachen übersetzt werden. Zudem besitzt Python einfach integrierte Funktionen für die Berechnung diskreter Exponentialfunktionen. Python ist zudem in der Lage vor der Kompilierung / Interpretierung und während der Laufzeit Funktionen zu Überschreiben. Python kann zudem objektorientiert verwendet werden, was bei der Abstraktion des Codes hilft und somit die Verständlichkeit fördert. Da der \ac{RSA}-Algorithmus mit großen ganzen Zahlen arbeiten muss bietet Python ein Vorteil im Gegensatz zu Alternativen wie Java, indem der zu Verfügung stehende Datentype int unbounded als von seiner Speicherlänge nicht begrenzt ist. In Java würde ein int maximal 32-bit Länge haben, was nicht ausreichend für jeglichen \ac{RSA}-Algorithmus ist. Zudem ist eine Deklaration von Datentypen in Python nicht nötig. 

        Als Angriffsziel wurde das PIP-Packet "rsa" gewählt. Dabei handelt es sich um eine Open Source Software Packet von Sybren A. Stüvel, welches den \ac{RSA}-Algorithmus und Hilfsfunktionalitäten komplett in Python implementiert. Dieses Packet wurde aufgrund der Beleibtheit ausgewählt. Es wird unter dem Namen "python-rsa" auf Git-hub gepflegt. Die Bearbeitung des Packet wurde mit dem Entwickler Sybren A. Stüvel abgesprochen.

        Nach Betrachtung des Aufbaus des Angriffsziels wurde sich dafür entschieden die Datei key.py manipulieren. Diese Python-Datei verwaltet die Generierung der Schlüssel(-parameter).
    
    \section{Angriffsvektoren}
        Die hier aufgeführten Angriffsvektoren beschreiben, wie ein kryptographisches System mittels dem hier entwickelten kleptographischen kompromittiert werden kann.

        \subsection{Malware}
            Malware sind bösartige Programme, welche Computer infizieren. Um den beschriebenen Angriff durchzuführen, muss die Malware in der Lage sein, lokale Dateien oder Pfade zu manipulieren. Ein möglicher Angriff kann wie folgt ablaufen:
            \begin{enumerate}
                \item Malware befällt das System
                \item Kleptographisch infiziertes Package wird heruntergeladen
                \item Das infizierte Package wird lokal mittel eines Package-Managers unter dem gleichen Namen wie das Original-Package installiert
                \item Dabei wird die Versionsnummer artifiziell hoch gesetzt, damit das Package immer bei "latest" genutzt wird
                \item Wenn lokale Programme nun das Package importieren, wird das kompromittierte Package verwendet
            \end{enumerate}
            Für diesen Angriff sind temporäre Schreibrechte auf dem System eine Mindestanforderung. Dadurch stellt sich die Frage, ob in diesem Fall andere Angriff auch bzw. besser geeignet sind. Ein Vorteil dieses Angriffes ergibt sich aus der Unübersichtlichkeit der von Dependency Chains. Zudem ist ein solcher Angriff vorteilhaft, wenn der Netzwerkverkehr des kompromittierten Systems überwacht wird. Eine solche Sicherheitsmaßnahme wäre keine Auswirkung, da die Rückgewinnung der Schlüsselparameter auf der selbstständigen Veröffentlichung des öffentlichen Schlüssels beruht.

        \subsection{Dependency Confusion}
            Als Dependency Confusion, werden Angriffe beschrieben, welche Versuchen durch Manipulation des Systems oder des Nutzers, eine bösartige Dependency in ein System einzuschleusen. Dies kann einerseits durch Veröffentlichung eines Packages mit ähnlichen Namen erfolgen oder durch das Vortäuschen neuerer Versionen. Hierfür könnte der Angreifer z.B. für das RSA Package von python "rsa":ein manipuliertes Package unter dem Namen "pyrsa" oder "rsa2". Dies könnte neue Nutzer dazu täuschen, das kompromittierte Package zu nutzen. Durch die Dependency Chain gilt das gleiche auch für Entwickler von Packages, welche eine RSA-Implementation benötigen. Falls es einem Angreifer gelingt, dass das kompromittierte Package eine Dependency eines anderen Packages wird, werden damit auch alle Packages und Nutzer kompromittiert, welche von dem Package mit der "schlechten" Dependency abhängig sind.

            Eine weitere Möglichkeit ist die Manipulation des Original-Packages. Dabei wird nicht ein eigenes Package erstellt, sondern versucht den Quellcode des echten Packages zu verändern. Dies kann einerseits durch das Hacken von Entwickleraccounts erfolgen aber auch durch das Beitragen zum Open Source Quellcode wobei die bösartigen Eigenschaften verborgen werden.
\chapter{Implementation}
    Es folgt die konkrete Implementation einer kleptographsichen Schwachstelle in die ausgewählte Softwarebibliothek.
\chapter{Risikoanalyse}
    In dieser Risikoanalyse (engl. Threat Assessment) wird die Gefahr analysiert und evaluiert, welche kleptographische Angriffe auf kryptographische Softwarebibliotheken bilden.

    Kleptographische Angriffe sind in der Lage, bei erfolgreichem Einsatz, kryptographische Systeme zo kompromittieren. Dadurch ist es für den Angreifer des kleptographischen Angriff möglich, die vom Nutzer verwendeten Krytposysteme zu brechen. Dadurch kann es zu Verletzung der Sicherheitsziele, Vertraulichkeit, Integrität und Authentizität, kommen. Ein Angreifer, kann somit in Namen des Opfers kommunizieren bzw. signieren, die Verschlüsslung zwischen Opfer und dritten Mitlesen und den Datenstrom manipulieren. Dies kann z.B. durch einen Man-In-The-Middle Angriff erfolgen wobei sich der Angreifer unentdeckt gegenüber zwei Kommunikationspartners als der jeweils andere ausgibt und als Zwischenstelle fungiert.

    Durch die \ac{SETUP} wird vermieden, dass versteckte Kanäle genutzt werden müssen. Dies ist einerseits hilfreich um Untersuchungen zu umgehen, welche z.B. Netzwerkverkehr auf diese analysieren. Andererseits weitet diese Eigenschaft den Einsatzbereich von kleptographischen Angriffen aus. Durch die Manipulation der Daten, welche selbstständig durch den Nutzer veröffentlicht werden, können auch Krytposysteme, welche keinen direkten Zugang zu versteckten Kanälen (hier: Netzwerk) haben, durch kleptographische Angriff kompromittiert werden. Dies ist besonders Sichtbar am Beispiel von \ac{TPM}s, da keine Netzwerkschnittstelle haben. Somit können durch kleptographische Angriffe sowohl stark überwachte als auch sehr isolierte Krytposysteme kompromittiert werden. 

    \section{Angriffsziele}
        \subsection{TPM}
            \ac{TPM} stellten die gefährlichsten Kryptosysteme dar, falls diese mit z.B. einer \ac{SETUP} versehen würde. \ac{TPM}s werden nur von wenigen Herstellern produziert und ihr Programmcode ist nach der Fertigung nicht manipulierbar. Zudem stellten \ac{TPM}s eine Back-Box dar, wodurch Quellcode-Analyse bzw. Reverse-Code-Engineering erschwert werden. Somit kann ein Angriff schlechter belegt bzw. erkannt werden. Eine kleptographische Backdoor, welche in einem \ac{TPM} entdeckt wird, impliziert mit hoher Wahrscheinlichkeit eine Kompromittierung aller \ac{TPM}s dieses Herstellers. Anderenfalls würde das Risiko bestehen, dass die Unterscheide zwischen \ac{TPM} über Seitenkanäle analysierbar sind.
            Der potenzielle Schaden durch eine Kompromittierung von \ac{TPM}s ist dabei am größten.

            Die Wahrscheinlichkeit, dass ein solcher Angriff erfolgreich wäre, ist sehr gering. Durch die Unveränderlichkeit des Programmcodes können solche Schwachstellen nicht nach der Herstellung erzeugt werden. Der Angriff müsste also vom Hersteller oder einem Teil dessen Lieferkette durchgeführt werden.

        \subsection{Krytpo-Bibliotheken}
            Krypto-Bibliotheken sind ein bisher wenig erforschtes Einsatzgebiet für kleptographische Angriffe. Dabei ist das Risiko zwischen proprietärer und Open-Source-Software zu unterscheiden.

            \paragraph{Proprietäre Software}
                Hersteller proprietärer Software können ähnlich zu Herstellern von \ac{TPM}s kleptographische Backdoor in ihre Produkte einbauen. Hierbei liegt auch eine Black-Box vor. Der Unterschied liegt in der Veränderbarkeit des Codes. Ein Hersteller oder ein Sourcecode-Virus kann nachträglich den Quellcode einer kleptographischen Backdoor versehen.
                Nutzer der Software können durch das geschlossene Design nur schwer überprüfen, ob sie angegriffen werden. Proprietäre Software bieten Herstellern dir Möglichkeit ihren Nutzern End-zu-End-Verschlüsselung zu versprechen und trotzdem Einsicht in die Kommunikation erlangen.

            \paragraph{Open-Source-Software}
                Open-Source-Software bietet Nutzern Sicherheit, indem sie in der Lage sind, selbst die verwendete Software auf Backdoors zu überprüfen. Jedoch ergibt sich durch die mehrstufigen Abhängigkeitsbeziehungen heutiger Open-Source-Software, die Möglichkeit eines Dependency Confusion Angriffs \ref{Dependency Confusion}.

    \section{Angriffsvektoren}
        Angriffsvektoren beschreiben die Wege über welche eine Angriff auf einem Zielsystem ausgeführt werden kann. 

        \subsection{Dependency Confusion} \label{Dependency Confusion}
            Grundfunktionen moderner Software jeglicher Hersteller basiert auf durch Open-Source-Software zur Verfügung gestellten Bibliotheken. Somit müssen sich Softwareentwickler nicht bei jedem Projekt neu damit befassen, Warteschlangen, Multithreading, Netzwerkprotokolle, kryptographische Funktionen, ... für ein neues Produkt zu entwickeln. Dieses Bereitstellen von Funktionalität wird durch Dependency (Abhängigkeiten) verwaltet. Dabei können Dependencies auch weitere Dependencies inherent haben. Somit spannt jede Dependency eines Softwareprodukts einen nicht zyklischen Graphen an Dependencies auf. Dabei werden Dependencies regelmäßig mit Sicherheitsupdates und neuen Features unter einer neuen Version aktualisiert. 
            Um auf dem neuesten Stand zu sein, wird von einer Dependency die neuste Version geladen.
            Dependency Confusion manifestiert sich auf zwei Arten. Der Entwickler des Produkts oder einer im Dependency Graphen des Produkts vorkommenden Dependencies verweist mit einer Dependency auf eine kompromittierte, andere Bibliothek. 
            Oder indem eine bestehenden Dependency im Graphen kompromittiert wird.

        \subsection{Sourcecode-Viren} \label{Sourcecode-Viren}
            Sourcecode-Viren verändern während ihrer Laufzeit den Quellcode einer Anwendung oder Dependency. Somit können lokale (zwischen-)gespeicherte Dependencies verändert bzw. kompromittiert werden. Der Angreifer braucht hierfür jedoch die Möglichkeit Code (den Virus) auf dem Zielsystem auszuführen. 



\chapter{Fazit}
    Die Sicherheit moderner kryptographischer Verfahren basiert auf der Geheimhaltung der Schlüssel und auf der Unberechenbarkeit schwerer mathematischer Probleme in vertretbarer Zeit.  
    Kleptographische Angriffe stellen ein Risiko für moderne Software dar. Durch einen \ac{SETUP} ist es Angreifern möglich geheime Daten zu erlangen, wobei nur die vom Nutzer selbst veröffentlichten Daten genutzt werden. 
    Kleptographie von kryptographischen Systemen ist dabei besonders gefährlich, da somit die Schutzziele der IT-Sicherheit, Vertraulichkeit, Integrität und Authentizität, verletzt werden können. Durch einen erfolgreichen Angriff auf ein RSA-Kryptosystem, kann der Angreifer sich als Opfer authentifizieren und Kommunikation entschlüsseln und manipulieren. 
    Der in dieser Arbeit behandelte, kleptographische Angriff erzeugt eine \ac{SETUP}. Dadurch wird der öffentliche Modulus eines erzeugten Schlüssels manipuliert, sodass die most significant Bits des Modulus dem verschlüsselten Primfaktor $P$ entsprechen. Somit veröffentlicht der Nutzer den manipulierten Schlüssel und somit den Primfaktor selbst, wenn er den öffentlichen Schlüssel publiziert. 
    Die Verschlüsselung ist einen asymmetrische RSA-Chiffre mit einem öffentlichen Schlüssel des Angreifers. Dadurch ist es nur dem Angreifer persönlich möglich den Primfaktor zu entschlüssel. Reverse-Engineering kann dabei nur den Angriff feststellen, aber nicht selbst benutzen.

    Der Angriff wurde erfolgreich in die Open-Source-Bibliothek "python-rsa" implementiert. Wenn ein Programm, welches in seinem Dependency Graphen auf diese kompromittierte Dependency verweist, einen Schlüssel erzeugt, kann ein Angreifer, wenn in Besitz des entsprechenden privaten Schlüssels die geheimen Schlüsselparameter bestimmen. Abgesehen von dieser Schwachstelle werden vollständig sichere und funktionierende \ac{RSA}-Schlüssel erzeugt.

    Die Algorithmen zur Erzeugung eines manipulierten Schlüssels und zur Extraktion der geheimen Parameter wurden im Laufe der Arbeit optimiert und genauer dokumentiert.

    Für einen erfolgreichen Angriff der hier aufgezeigten Art, benötigt eine Angreifer die Möglichkeit Code auf die Zielsystem auszuführen und damit Programmcode zu ändern oder die Manipulation der kryptographischen Softwarebibliothek oder von Dependencies, welche zur Schlüsselerzeugung auf dem Zielsystem genutzt wird.  

% Ab hier beginnt der Anhang
\appendix
\addcontentsline{toc}{chapter}{Anhang}

\addcontentsline{toc}{chapter}{Index}
\printindex

\addcontentsline{toc}{chapter}{Literaturverzeichnis}




% Haben Sie das "biblatex"-Paket nicht installiert, benutzen Sie folgendes:
% Ohne das "biblatex"-Paket (s. bericht.sty) produziert folgendes
% "deutsche" Zitate in Literaturverzeichnissen gemaß der Norm DIN 1505,
% Teil 2 vom Jan. 1984.
% Die Zitatmarken werden alphabetisch nach Verfassern
% sortiert und sind durch abgekürzte Verfasserbuchstaben plus
% Erscheinungsjahr in eckigen Klammern gekennzeichnet.

% \bibliographystyle{alpha}
% \bibliography{bericht}

%%%%%%%%%%%%%%%%%%%%%%%%%%%%%%%%%%%%%%%5
% BIBLATEX
% Benutzt man das "biblatex"-Paket, muß man folgendes schreiben:
\def\refname{Literaturverzeichnis}
\printbibliography
%%%%%%%%%%%%%%%%%%%%%%%%%%%%%%%%%%%%%%%5

\end{document}
